\subsection{Family monodromy and occupancy}

\begin{definition}[Family space]
Let $\overline X_f$ denote the compact seam quotient with four $\mathbb Z_4$ corner points
$p_1,\dots,p_4$, and define the punctured family section
\[
X_f^\circ:=\overline X_f\setminus\{p_1,p_2,p_3,p_4\}.
\]
Its rigid harmonic model is the Möbius class
\[
X_f^\circ\cong \mathbb P^1\setminus\mu_4.
\]
Set
\[
F:=H^1(X_f^\circ,\mathbb C).
\]
We use the distinguished cycle basis $(C_1,C_2,C_T)$ only as a concrete carrier-side basis of this family space; any later APS/Hodge realization is read as compatible analytic structure on the same topological family space.
\end{definition}
In this interpretation $F$ is neither an arbitrary flavor vector space nor a loose bookkeeping
trick; it is a cycle space attached to the seam data. The actual bridge logic is shown in
\Cref{fig:family-triangle}: the derived corner theorem fixes $N_{\mathrm{fam}}=3$, the
topological family step promotes that count to a family space, and the occupancy corollary then
feeds back into the defect channel.

\begin{figure}[t]
\centering
\scriptsize
\begin{tikzpicture}[
    >=Latex,
    hard/.style={draw=blue!60!black, fill=blue!6, rounded corners, align=center, text width=2.2cm, minimum height=1.0cm},
    closure/.style={draw=orange!75!black, fill=orange!8, rounded corners, align=center, text width=2.4cm, minimum height=1.0cm},
    bridge/.style={draw=green!50!black, fill=green!8, rounded corners, align=center, text width=2.25cm, minimum height=1.0cm}
]
\draw[draw=blue!60!black, fill=blue!4, rounded corners] (-5.5,-1.4) rectangle (-3.0,1.4);
\foreach \x/\y in {-5.1/1.0,-3.4/1.0,-5.1/-1.0,-3.4/-1.0} {
    \filldraw[blue!60!black] (\x,\y) circle (2.3pt);
    \node[blue!60!black, anchor=south] at (\x,\y) {\tiny $\frac34$};
}
\node[align=center] at (-4.25,0) {\textbf{corner compactification}\\four $\mathbb Z_4$ corners\\each contributes $\frac34$};

\node[hard] (nfam) at (-1.0,0) {$N_{\mathrm{fam}}$\\$=4\cdot\frac34=3$};

\coordinate (A) at (2.0,1.15);
\coordinate (B) at (0.7,-0.95);
\coordinate (C) at (3.3,-0.95);
\node[closure, text width=1.2cm, minimum height=0.8cm] at (A) {$C_1$};
\node[closure, text width=1.2cm, minimum height=0.8cm] at (B) {$C_2$};
\node[closure, text width=1.2cm, minimum height=0.8cm] at (C) {$C_T$};
\draw[->, thick, orange!75!black] ($(A)+(0.2,-0.3)$) .. controls (2.7,0.25) .. ($(C)+(-0.2,0.3)$);
\draw[->, thick, orange!75!black] ($(C)+(-0.38,-0.02)$) .. controls (2.0,-1.65) .. ($(B)+(0.38,-0.02)$);
\draw[->, thick, orange!75!black] ($(B)+(0.2,0.3)$) .. controls (1.3,0.2) .. ($(A)+(-0.2,-0.3)$);
\node[closure, text width=2.1cm, minimum height=0.9cm] at (2.0,-2.2) {topological family closure\\$F\cong\mathbb C^3$};

\node[closure] (occ) at (5.85,0.55) {$\Omega_{\mathrm{adm}}$\\$=3\times16=48$};
\node[bridge] (dtop) at (5.85,-1.15) {$\deltatop=\Omega_{\mathrm{adm}}\cthree^4$\\boundary-to-bulk bridge};

\draw[->, thick] (-3.0,0) -- (nfam);
\draw[->, thick] (nfam) -- (1.0,0.1);
\draw[->, thick] (nfam) -- (occ);
\draw[->, thick, orange!75!black] (2.0,-1.55) -- (2.0,-1.8);
\draw[->, thick, orange!75!black] (3.05,-1.9) -- (occ);
\draw[->, thick, green!50!black] (occ) -- (dtop);
\end{tikzpicture}
\caption{Carrier-to-family closure bridge. Corner compactification yields the punctured family
section $X_f^\circ$, the harmonic family space $\mathcal H_F\cong F\cong H^1(X_f^\circ,\mathbb C)$
gives the hard multiplicity $N_{\mathrm{fam}}=3$, and the resulting occupancy
$\Omega_{\mathrm{adm}}=48$ feeds back into the bridge reading
$\deltatop=\Omega_{\mathrm{adm}}\cthree^4$.}
\label{fig:family-triangle}
\end{figure}

The Hodge metric on $F$ turns the monodromy data into an honest unitary family system rather than a merely decorative triangle sketch.

\begin{theorem}[Four corners and the rigid family surface \texorpdfstring{$\mathbb P^1\setminus\mu_4$}{P1 minus mu4}]
\label{thm:topological-family-closure}
Let $\overline X_f$ be the compact seam quotient induced by the deck involution $\taudbl$ and its
spin lift on the oriented normal two-plane. Every admissible fixed corner contributes one
quarter-turn of the lifted spin rotation. Therefore
\[
N_{\mathrm{corner}}\cdot \frac{\pi}{2}=2\pi,
\qquad
N_{\mathrm{corner}}=4.
\]
Removing the four corner points gives a four-punctured sphere
\[
X_f^\circ:=\overline X_f\setminus\{p_1,p_2,p_3,p_4\}.
\]
The inherited finite symmetry group on the puncture set is a faithful dihedral group
$D_4\subset PGL_2(\mathbb C)$. Every faithful copy of $D_4$ in $PGL_2(\mathbb C)$ is M\"obius
conjugate to the standard action generated by
\[
r(z)=iz,
\qquad
s(z)=1/z.
\]
Up to M\"obius equivalence the unique $D_4$-invariant four-point configuration is therefore
\[
\mu_4:=\{z\in\mathbb C:\ z^4=1\}=\{1,-1,i,-i\}.
\]
Hence
\[
X_f^\circ\cong \mathbb P^1\setminus\mu_4,
\qquad
\chi(X_f^\circ)=-2,
\qquad
H_1(X_f^\circ,\mathbb Z)\cong \mathbb Z^3.
\]
Consequently
\[
F:=H^1(X_f^\circ,\mathbb C)\cong \mathbb C^3.
\]
\end{theorem}

\begin{proof}
The quarter-turn count gives four corner points. A sphere with four removed points is isomorphic to
$\mathbb P^1\setminus\{q_1,q_2,q_3,q_4\}$ for some ordered quadruple. The classification of finite
subgroups of $PGL_2(\mathbb C)$ implies that every faithful dihedral subgroup of order eight is
M\"obius conjugate to the standard $D_4$ action generated by $z\mapsto iz$ and $z\mapsto 1/z$.
Under that action the unique invariant orbit of cardinality four is the set of fourth roots of
unity, namely $\mu_4=\{1,-1,i,-i\}$. Therefore
\[
X_f^\circ\cong \mathbb P^1\setminus\mu_4.
\]
The Euler characteristic is $2-4=-2$, and for a genus-zero surface with four punctures one has
$\rank H_1=4-1=3$. Dualizing gives $F\cong\mathbb C^3$.
\end{proof}

\begin{theorem}[Three harmonic family modes on the four-punctured family surface]
\label{thm:harmonic-family-modes}
Let
\[
X_f^\circ=\mathbb P^1\setminus\mu_4
\]
with a complete cusp metric in its M\"obius class, and define the family Hilbert space by
\[
\mathcal H_F:=\ker \Delta^{(1)}_{L^2,\mathrm{APS}}(X_f^\circ).
\]
Then
\[
\dim\mathcal H_F=3.
\]
If the physical family multiplicity is defined by
\[
N_{\mathrm{fam}}:=\dim\mathcal H_F,
\]
then
\[
N_{\mathrm{fam}}=3,
\qquad
F\cong \mathcal H_F\cong \mathbb C^3.
\]
\end{theorem}

\begin{proof}
For a compact truncation $\overline X_{f,R}$ with four boundary circles,
\[
\chi(\overline X_{f,R})=2-4=-2.
\]
Hence
\[
\dim H^1(\overline X_{f,R},\partial\overline X_{f,R};\mathbb C)=3.
\]
The $L^2$ Hodge theorem for finite-area cusp surfaces identifies
\[
\ker \Delta^{(1)}_{L^2,\mathrm{APS}}(X_f^\circ)
\cong
H^1(\overline X_{f,R},\partial\overline X_{f,R};\mathbb C),
\]
so the dimension is $3$.
\end{proof}

\begin{remark}[How the family theorem is now used]
The family theorem is used at rigid conformal strength in the present version. Any orbifold APS
realization is read as compatible analytic structure on the same classification-rigid
four-punctured section, while the physical multiplicity itself is carried by the harmonic family
mode space $\mathcal H_F$.
\end{remark}

\begin{proposition}[Orbifold APS compatibility]
\label{prop:orbifold-aps-compatibility}
Under the additional analytical hypotheses of \TFPTxref{app:aps},
\[
\Ind^{\mathrm{orb,APS}}(D_{\mathrm{rel}}^+\otimes L_F)=3.
\]
\end{proposition}

\begin{corollary}[Boundary-to-bulk bridge corollary for the admissible occupancy]
\label{cor:boundary-bulk-occupancy}
\label{cor:derived-occupancy-identity}
On the canonical branch,
\[
\dim(S^+\otimes \mathcal H_F)=16\cdot 3=48=\Oadm.
\]
Hence the topological surcharge enters the present paper as a derived occupancy corollary:
\[
\deltatop
=
\cthree^4\,\dim(S^+\otimes \mathcal H_F)
=
\Oadm\,\cthree^4.
\]
\end{corollary}

The family subsection is therefore no longer split between a local count, an independent index
theorem, and a later occupancy reinterpretation. Those statements are compressed into one derived
corner theorem with one bridge corollary, while the holonomy proposition below records the
resulting unitary geometry on $F\cong \mathcal H_F$.

This theorem is the family statement referenced later by the admissibility theorems: the closure layer uses it as the explicit family component of the master closure theorem rather than as an external import.

\begin{proposition}[Rigid special-unitary family holonomy on the determinant-trivial branch]
\label{prop:family-holonomy-su3}
Let $L_F\to X_f^\circ$ be the flat Hermitian family bundle with fiber
$F\cong \mathcal H_F\cong H^1(X_f^\circ,\mathbb C)$. Assume that its determinant line is flat trivial,
\[
\det L_F \cong X_f^\circ\times \mathbb C,
\]
and that the chosen flat connection preserves this trivialization. Then
\[
\mathrm{Hol}(\nabla_F)\subset SU(3)_F.
\]
\end{proposition}

\begin{proof}
The Hermitian metric implies
\[
\mathrm{Hol}(\nabla_F)\subset U(3).
\]
The induced connection on $\det L_F$ has trivial holonomy by assumption, so for
every loop $\gamma$ one has
\[
\det \rho_F(\gamma)=1.
\]
Hence every holonomy matrix lies in $SU(3)$.
\end{proof}

The role of the present subsection is thus no longer only mnemonic. It displays both the local patch-deficit count and the global APS/Hodge realization before the later closure stack compresses them into the theorem-level output surface.

\subsection{Occupancy and boundary-to-bulk bridge}

With the family bridge installed, the matter space becomes $S^+\otimes \mathcal H_F$ and the
occupancy $\Oadm=48$ is already fixed by the harmonic family-mode theorem and its bridge
corollary. Here \emph{occupancy} means the admissible chiral-state count after family lifting; it
is not used as a primitive datum. The bridge point is not that family lifting creates a new
numerical surcharge. Rather, the inherited hard kernel value
\[
\deltatop=\frac{3}{256\pi^4}
\]
admits on the canonical family branch the occupancy reinterpretation $\deltatop=\Oadm\,\cthree^4$.
Likewise the hard second-order defect contribution is re-read as
\[
\delta_2=\frac54\,\deltatop^2=B\gamma\,\Oadm^2\cthree^8.
\]
In the present manuscript this occupancy corollary is used with a derived bridge reading:
it is the boundary-to-bulk step that later feeds both the electromagnetic closure and the
geometric channel. The coefficient $5/4$ is therefore not an extra numerical motif; it is the
minimal-carrier compression factor $B\gamma$ rewritten in defect language. At the same time, the
equality $\deltatop=\Oadm\cthree^4$ is not presented as an independent dynamical transgression
theorem in the ordinary field-theory sense.

\subsection{From bridge to geometric completion}

\begin{principle}[Closure split]
Once the carrier packet and family closure architecture are fixed, the focused main-text closure
problem splits into a geometric program $\Ageom$ and a transport program $\Atrans$. Record
algebra and prediction semantics are now tied directly to the admissible gap in the main theorem
stack, while E$_8$ scale grammar and horizon extrapolations are deferred to appendix-level
continuations.
\end{principle}

\[
D_{\mathrm{geo}}:=D_{\mathrm{ref}}+D_{\mathrm{rel}}=D.
\]
The two main-text programs are packaged as
\[
\Ageom
:=
\langle E_3\oplus E_2,\ D_{\mathrm{geo}},\ \mathcal B_{\mathrm{rel}},\ \wspin\rangle,
\qquad
\Atrans
:=
\langle U_6,\ T,\ D_y,\ \mathcal Y_y^{(\varepsilon)},\ \Ccyc\rangle.
\]
This split keeps the carrier theorem, the geometric reconstruction problem, and the record-level extensions from being flattened into one undifferentiated claim. In particular, the metric field is no longer treated here as an independent entry of $\Ageom$; it is reconstructed from the geometric Dirac anchor introduced next, while $D_{\mathrm{rel}}$ remains the comparison operator for spectral subtraction.

\subsection{Relative APS and superconnection setup}

The geometric completion is formulated on a graded bundle
\[
\mathcal E_{\mathrm{rel}}=\mathcal E_{\mathrm{rel}}^+\oplus\mathcal E_{\mathrm{rel}}^-,
\]
with the geometric Dirac anchor $D_{\mathrm{geo}}$, the declared reference operator
$D_{\mathrm{ref}}$, and APS-type boundary conditions on the seam $\Seam$; see
Appendix~\TFPTxref{app:aps}. The geometric package is
\[
\mathbb A_{\Seam}:=(D_{\mathrm{geo}},D_{\mathrm{rel}},A_{\Sigma}^{\mathrm{geo}}),
\qquad
A_{\Sigma}^{\mathrm{geo}}:=\mathcal B_{\mathrm{rel}}\oplus\nabla_{\chigeo}.
\]
Here $D_{\mathrm{geo}}$ is the geometric Dirac anchor, while $D_{\mathrm{rel}}$ is the relative comparison operator entering the spectral subtraction formulas. The bundle data $A_{\Sigma}^{\mathrm{geo}}$ reorganize the primitive seed package $A_{\Sigma}^{\mathrm{seed}}$ introduced earlier in \TFPTxref{eq:tfpt-master}.
\begin{theorem}[Geometric reconstruction from the Dirac anchor]
\label{thm:geometric-reconstruction-anchor}
\label{thm:geometric-reconstruction-from-Dgeo}
\label{thm:relative-geometric-reconstruction}
Assume
\[
Z(\mathcal A)\cong C^\infty(M),
\]
and the geometric anchor $D_{\mathrm{geo}}$ is of Dirac type with principal symbol
\[
\sigma(D_{\mathrm{geo}})(x,\xi)^2
=
g^{\mu\nu}(x)\xi_\mu\xi_\nu\,\mathbf 1.
\]
Then $D_{\mathrm{geo}}$ reconstructs the metric $g$, the spin structure, and the
Levi--Civita connection on the geometric branch. Inner fluctuations of the same anchored
spectral datum generate the gauge sector.
The operator $D_{\mathrm{rel}}$ is the reference-subtracted comparison operator on that
same branch and is not itself used as the local geometric anchor.
\end{theorem}

\begin{proof}[Proof sketch]
On the admissible branch the commutative center identifies the manifold algebra, while the
principal symbol of the Dirac anchor determines the quadratic form on cotangent vectors.
Standard Dirac-type reconstruction then recovers the spin bundle and Levi--Civita connection
from that same symbol data. The relative operator changes the comparison bookkeeping, not the
geometric source of the local symbol.
\end{proof}

\begin{theorem}[Retained branch from the admissibility selector]
\label{thm:retained-branch-from-Padm}
\label{thm:admissible-sector-on-reconstructed-geometry}
\label{thm:full-operator-reconstruction}
Let $D_{\mathrm{geo}}$ be the geometric anchor and let $\Padm$ be the retained
admissibility selector.
Then:
\begin{enumerate}[label=(\roman*), leftmargin=1.5em, topsep=3pt, itemsep=2pt]
\item $\Xbulk=\mathrm{Spec}\,Z(\mathcal A)$ and $\Seam=\mathrm{Fix}(\taudbl)$;
\item the local spectral density of $D_{\mathrm{geo}}$ determines the geometric response
      $\chigeo$;
\item $\Padm$ defines the retained admissible sector for observables, transport,
      and positivity statements on that branch.
\end{enumerate}
No claim is made that the compressed operator $\Padm D_{\mathrm{rel}}\Padm$
carries the local principal symbol of the geometry.
\end{theorem}

\begin{corollary}[No residual geometric datum beyond the geometric anchor]
The geometric package is reconstructed from $D_{\mathrm{geo}}$ and its local spectral density.
The selector $\Padm$ acts downstream on the retained sector and is not part of the
local geometric anchor.
\end{corollary}

\begin{definition}[Local spectral scale and canonical relative spectral residues]
\label{def:local-spectral-scale}
\label{def:canonical-residues}
Write the geometric response as a local spectral scale
\[
\chigeo=\Lambda e^\sigma,
\qquad
D_{\sigma,\mathrm{geo}}:=e^{-\sigma/2}D_{\mathrm{geo}}e^{-\sigma/2}.
\]
\[
\mathcal Z_{\mathrm{rel}}(s;\sigma)
:=
\operatorname{Tr}_{\mathrm{rel}}\!\left(
|D_{\sigma,\mathrm{geo}}+\mathcal B_{\mathrm{rel}}|^{-2s}
-|D_{\mathrm{ref}}|^{-2s}
\right)
+\frac{i\pi}{2}\Delta\eta_\Sigma.
\]
The canonical Einstein residue and order-zero residue are
\[
\mathcal E_D
:=
\operatorname*{Res}_{s=1}\mathcal Z_{\mathrm{rel}}(s;\sigma),
\qquad
\Gamma_D^{(4)}
:=
\operatorname{FP}_{s=0}\mathcal Z_{\mathrm{rel}}(s;\sigma).
\]
The gravitational branch is then read from the profile-free residue pair rather than from a freely
chosen spectral profile.
\end{definition}

\begin{corollary}[Constant-scale reduction of the geometric branch]
\label{cor:constant-scale-reduction}
If the local spectral scale is frozen to a constant $\sigma=\sigma_0$, then
$\chigeo=\chigeozero:=\Lambda e^{\sigma_0}$ and the present geometric branch reduces to the
constant-$\chigeo$ closure relations extracted from
$\Gamma_{\mathrm{grav}}:=-6\chigeo^2\mathcal E_D+\Gamma_D^{(4)}$:
\[
\Mpl^2=\frac{\chigeozero^2}{2\pi^2},
\qquad
\vgeo^2=Z_H\,\frac{\mu_\Phi^2(\chigeozero)}{\lambda_\Phi}.
\]
Once the unique stationary root of \TFPTxref{thm:absolute-spectral-planck-closure} is imposed, the
Einstein coefficient is equivalently the boundary-normalized readout
\[
\frac{\Mpl^2}{\lambda_\Sigma^2}=\frac{\rho_\star}{2\pi^2}.
\]
Thus the former constant-$\chigeo$ formulas are retained as the homogeneous branch of the local
spectral-scale formulation rather than as a competing package.
\end{corollary}


